\documentclass[11pt,a4paper]{article}
\newcommand{\Rodzaj}{Karta pomiarowa}
\usepackage{fontspec}
\setmainfont{texgyretermes-regular.otf}[BoldFont=texgyretermes-bold.otf,ItalicFont=texgyretermes-italic.otf]
\usepackage[polish]{babel}
\usepackage{amsmath,amssymb,geometry,array,tabularx,graphicx,caption,fancyhdr}
\usepackage[table]{xcolor}
\geometry{left=18mm,right=18mm,top=19mm,bottom=19mm,headheight=14pt,headsep=6mm}
\setlength{\parindent}{0pt}\setlength{\parskip}{5pt}
\setlength{\emergencystretch}{2em}
\renewcommand{\arraystretch}{1.3}\setlength{\tabcolsep}{5pt}
\setlength{\fboxsep}{6pt}
\setlength{\abovedisplayskip}{8pt}\setlength{\belowdisplayskip}{8pt}
\setlength{\abovedisplayshortskip}{5pt}\setlength{\belowdisplayshortskip}{5pt}
\widowpenalty=10000\clubpenalty=10000
\raggedbottom
\newcolumntype{Y}{>{\raggedright\arraybackslash}X}
\newcolumntype{C}{>{\centering\arraybackslash}X}
\newcommand{\naglowek}{\rowcolor{gray!12}}
\newcommand{\odstep}{\par\smallskip}
\captionsetup{font=small,labelfont=bf,justification=justified,singlelinecheck=false}
\pagestyle{fancy}\fancyhf{}\fancyhead[L]{\small Ćwiczenie 83 — spektrometr optyczny}
\fancyhead[R]{\small\Rodzaj}\fancyfoot[C]{\small\thepage}
\renewcommand{\headrulewidth}{0.3pt}
\makeatletter
\renewcommand\section{\@startsection{section}{1}{0pt}{10pt}{4pt}{\large\bfseries}}
\renewcommand\subsection{\@startsection{subsection}{2}{0pt}{8pt}{4pt}{\normalsize\bfseries}}
\makeatother
\newcommand{\pole}{\rule{23mm}{0.3pt}}
\newcommand{\miejsce}[2]{\par\medskip\noindent\fbox{\begin{minipage}[t][#1][t]{\dimexpr\linewidth-2\fboxsep-2\fboxrule\relax}\small #2\end{minipage}}\par\medskip}
\newcommand{\metryka}{%
\begingroup\small\setlength{\parskip}{0pt}\renewcommand{\arraystretch}{1.1}
\begin{tabularx}{\linewidth}{|p{29mm}|Y|p{18mm}|p{18mm}|p{19mm}|}\hline
\textbf{Wydział}\newline Informatyki & \textbf{Imię i nazwisko}\newline 1. Marcin Podyma\newline 2. Paweł Mierzwa & \textbf{Rok}\newline 2026 & \textbf{Grupa}\newline 3 & \textbf{Zespół}\newline 1\\\hline
\end{tabularx}\par\nointerlineskip
\begin{tabularx}{\linewidth}{|p{29mm}|Y|p{19mm}|}
\parbox[t]{29mm}{\centering\textbf{PRACOWNIA\\FIZYCZNA\\WFiIS AGH}} & \textbf{Temat:}\newline Spektrometr optyczny — widmo liniowe par rtęci & \textbf{Nr}\newline\textbf{ćwiczenia}\newline 83\\\hline
\end{tabularx}\par\nointerlineskip
\begin{tabularx}{\linewidth}{|p{29mm}|Y|Y|Y|Y|p{19mm}|}
\textbf{Data wykonania}\newline\strut & \textbf{Data oddania}\newline\strut & \textbf{Zwrot do popr.} & \textbf{Data oddania} & \textbf{Data zaliczenia} & \textbf{OCENA}\\\hline
\end{tabularx}\endgroup\medskip}

\usepackage{multirow}
\newcommand{\notatki}[1]{%
\par\smallskip\textbf{#1}\par
\rule{\linewidth}{0.25pt}\par\vspace{5mm}
\rule{\linewidth}{0.25pt}\par\vspace{5mm}
\rule{\linewidth}{0.25pt}\par}
\newcommand{\zeroserii}{%
\begin{tabularx}{\linewidth}{|Y|p{27mm}|p{27mm}|}\hline
\naglowek\textbf{Wiązka nieodchylona — odczyt $\theta_0$} & \textbf{Stopnie [$^\circ$]} & \textbf{Minuty [$'$]}\\\hline
Kierunek zerowy dla tej serii & &\rule{0pt}{9mm}\\\hline
\end{tabularx}\par\medskip}
\newcommand{\tabelaodczytow}{%
\begingroup\small\setlength{\tabcolsep}{5pt}\renewcommand{\arraystretch}{1}
\begin{tabularx}{\linewidth}{|c|p{35mm}|c|C|C|C|C|}\hline
\naglowek\textbf{ID} & \textbf{Barwa linii} & \textbf{Pomiar} &
\multicolumn{2}{c|}{\textbf{Lewa strona $\theta_L$}} &
\multicolumn{2}{c|}{\textbf{Prawa strona $\theta_P$}}\\\cline{4-7}
 & & & stopnie [$^\circ$] & minuty [$'$] & stopnie [$^\circ$] & minuty [$'$]\\\hline
\csname @@input\endcsname{odczyty.tex}
\end{tabularx}\endgroup\par}
\begin{document}
\metryka
\section*{1. Parametry stanowiska}
\begin{tabularx}{\linewidth}{|Y|p{53mm}|p{20mm}|}\hline
\naglowek\textbf{Parametr} & \textbf{Wartość} & \textbf{Jednostka}\\\hline
Liczba rys siatki $N$ & & $\mathrm{mm}^{-1}$\rule{0pt}{10mm}\\\hline
Stała siatki $d=1/N$ & & mm\rule{0pt}{10mm}\\\hline
Szerokość szczeliny & & mm\rule{0pt}{10mm}\\\hline
Przyjęta dokładność odczytu kąta & & minuty ($'$)\rule{0pt}{10mm}\\\hline
\end{tabularx}
\odstep\small Stałą siatki lub liczbę rys odczytaj na stanowisku.
Nie wpisuj parametrów z przykładowego sprawozdania bez sprawdzenia.\normalsize

\section*{2. Kontrola ustawienia}
Dla tej samej linii w I rzędzie porównaj odchylenia lewej i prawej
strony względem wiązki nieodchylonej. Różnica większa niż około
$10'$ wymaga poprawienia prostopadłości siatki.
\par\medskip
\begin{tabularx}{\linewidth}{|Y|C|C|C|}\hline
\naglowek\textbf{Linia kontrolna / barwa} & $\alpha_L$ [$^\circ$] & $\alpha_P$ [$^\circ$] & \textbf{Różnica [$'$]}\\\hline
 & & &\rule{0pt}{13mm}\\\hline
\end{tabularx}
\odstep
$\square$ ostrość lunety\quad $\square$ odczyt zera\quad
$\square$ prostopadłość siatki
\notatki{Uwagi do przygotowania stanowiska}

\section*{3. Zasady zapisu}
\begin{itemize}\setlength{\itemsep}{2pt}
\item W tabelach zapisuj \textbf{surowe wskazania skali $\theta$}.
Odchylenia $\alpha$ obliczysz później względem $\theta_0$.
\item Stopnie i minuty wpisuj w osobnych kolumnach.
Przykład: $14^\circ20'$ to $14{,}3333\ldots{}^\circ$, nie $14{,}20^\circ$.
\item Zachowaj ten sam ID dla danej linii po obu stronach i w obu rzędach.
Obie linie żółte wpisz jako dwa osobne ID.
\item Trzy wiersze pod jednym ID służą powtórzeniom. Niewykonane
pomiary pozostaw puste, a niewidoczną linię oznacz w uwagach.
\end{itemize}

\clearpage
\section*{I rząd widma — $n=1$}
\zeroserii
\small Wpisuj wskazania skali. Dwa żółte prążki mają osobne ID.
Trzy pomiary w grupie dotyczą tej samej linii.\normalsize
\par\medskip\tabelaodczytow
\odstep\small Uwagi: \rule{\dimexpr\linewidth-14mm\relax}{0.25pt}\normalsize

\clearpage
\section*{II rząd widma — $n=2$}
\zeroserii
\small Zachowaj ID z I rzędu. Jeśli linii nie widać, wpisz
„niewidoczna” w polu barwy i nie uzupełniaj odczytu.\normalsize
\par\medskip\tabelaodczytow
\odstep\small Uwagi: \rule{\dimexpr\linewidth-14mm\relax}{0.25pt}\normalsize
\end{document}
